\chapter{Rekursion und Charakterisierung}
\label{b10guide:ch:recursion}

\section{Rekursion auf den natürlichen Zahlen}

Die Zielmengen in diesem Kapitel tragen keine zusätzliche axiomatische
Struktur. Anfangswerte und Schrittvorschriften sind jeweils ausdrücklich
angegebene Daten. Die Sätze betreffen die durch natürliche Zahlen
festgelegte Reihenfolge der Definition. Die klassische Grundlage
für Rekursion und Eindeutigkeit findet sich bei Dedekind \cite{dedekind}.

\begin{theorem}[Dedekindscher Rekursionssatz]
\label{b10guide:thm:dedekind}
Seien $A$ eine Menge, $a\in A$ und $F\colon A\to A$.
Dann existiert genau eine Funktion $u\colon\mathbb N\to A$ mit
\[
 u(0)=a,\qquad u(S(n))=F(u(n))\quad(n\in\mathbb N).
\]
\end{theorem}

\begin{theorem}[Rekursion mit vom Schritt abhängiger Vorschrift]
\label{b10guide:thm:step-recursion}
Seien $A$ eine Menge, $a\in A$ und
$H\colon\mathbb N\times A\to A$.
Dann existiert genau eine Funktion $u\colon\mathbb N\to A$ mit
\[
 u(0)=a,\qquad u(S(n))=H(n,u(n))\quad(n\in\mathbb N).
\]
\end{theorem}

\begin{theorem}[Rekursion mit Parametern]
\label{b10guide:thm:parameter-recursion}
Seien $P,A$ Mengen, $a\colon P\to A$ und
$H\colon P\times\mathbb N\times A\to A$.
Dann existiert genau eine Funktion $u\colon P\times\mathbb N\to A$
mit
\[
 \begin{aligned}
 u(p,0)&=a(p),\\
 u(p,S(n))&=H(p,n,u(p,n))
 \end{aligned}
 \qquad(p\in P,\ n\in\mathbb N).
\]
\end{theorem}

\begin{theorem}[Simultane Rekursion zweier Größen]
\label{b10guide:thm:simultaneous-recursion}
Seien $a\in A$, $b\in B$ und
\[
 H\colon\mathbb N\times A\times B\to A,
 \qquad K\colon\mathbb N\times A\times B\to B.
\]
Dann existiert genau ein Paar von Funktionen
$u\colon\mathbb N\to A$, $v\colon\mathbb N\to B$ mit
\[
 \begin{aligned}
 u(0)&=a,\qquad v(0)=b,\\
 u(S(n))&=H(n,u(n),v(n)),\\
 v(S(n))&=K(n,u(n),v(n))
 \end{aligned}
 \quad(n\in\mathbb N).
\]
\end{theorem}

\begin{definition}[Anfangsabschnitt und bisherige Werte]
\label{b10guide:def:recursion-history}
Für $n\in\mathbb N$ sei
\[
 I_n=\{k\in\mathbb N\mid k<n\}.
\]
Für eine Menge $A$ sei
\[
 \mathcal H_A=\{(n,h)\mid n\in\mathbb N,\ h\colon I_n\to A\}.
\]
Für $u\colon\mathbb N\to A$ bezeichnet $u|_{I_n}$ die Funktion
auf $I_n$, die jedem $k\in I_n$ den Wert $u(k)$ zuordnet.
Sie enthält genau die vor dem Schritt $n$ bestimmten Werte.
\end{definition}

\begin{theorem}[Rekursion unter Verwendung aller bisherigen Werte]
\label{b10guide:thm:history-recursion}
Seien $A$ eine Menge und $\Phi\colon\mathcal H_A\to A$.
Dann existiert genau eine Funktion $u\colon\mathbb N\to A$ mit
\[
 u(n)=\Phi(n,u|_{I_n})\quad(n\in\mathbb N).
\]
Insbesondere ist $u(0)=\Phi(0,\varnothing)$.
\end{theorem}

\begin{theorem}[Erhaltung einer Eigenschaft durch Rekursion]
\label{b10guide:thm:recursion-invariant}
Seien $u\colon\mathbb N\to A$, $a\in A$ und
$H\colon\mathbb N\times A\to A$ mit
$u(0)=a$ und $u(S(n))=H(n,u(n))$ für alle $n\in\mathbb N$.
Für $B\subseteq A$ gelte
\[
 a\in B,\qquad
 \forall n\in\mathbb N\;\forall x\in B\;(H(n,x)\in B).
\]
Dann gilt $u(n)\in B$ für alle $n\in\mathbb N$.
\end{theorem}

\section{Eindeutigkeit der natürlichen Zahlen}

\begin{theorem}[Eindeutiger Vergleich zweier Peano-Systeme]
\label{b10guide:thm:categoricity}
Seien $(N,0_N,S_N)$ und $(M,0_M,S_M)$ Peano-Systeme im Sinn
von Definition~\ref{b10guide:def:peano}.
Dann existiert genau eine Funktion $f\colon N\to M$ mit
\[
 f(0_N)=0_M,\qquad f(S_N(n))=S_M(f(n))\quad(n\in N).
\]
Für diese Funktion besitzt jedes $m\in M$ genau ein Urbild in $N$.
Somit bestimmen die Peano-Eigenschaften die natürlichen Zahlen bis
auf eine eindeutige Umbenennung, die Null und Nachfolger erhält.
\end{theorem}

\begin{theorem}[Keine nichttriviale Umbenennung von Null und Nachfolger]
\label{b10guide:thm:rigidity}
Gilt für $f\colon\mathbb N\to\mathbb N$
\[
 f(0)=0,\qquad f(S(n))=S(f(n))\quad(n\in\mathbb N),
\]
so gilt $f(n)=n$ für alle $n\in\mathbb N$.
\end{theorem}

